\setauthor{Elisa Kaltenberger}

Die Projektidee entstand aus eigenem Interesse und der grundlegenden Fragestellung, wie Barrierefreiheit im Alltag tatsächlich erlebt wird, insbesondere aus der Perspektive von Rollstuhlfahrenden. Im schulischen und gesellschaftlichen Kontext wird Barrierefreiheit häufig theoretisch behandelt, jedoch selten praktisch erfahrbar gemacht.

Ein weiterer Ausgangspunkt war das Interesse an immersiven Technologien wie Virtual Reality. VR bietet die Möglichkeit, Perspektiven realitätsnah zu simulieren und Nutzer*innen unmittelbar in eine andere Lebenssituation zu versetzen. Dadurch entstand die zentrale Leitfrage:

\textbf{Wie kann man typische Hindernisse für Rollstuhlfahrende so darstellen, dass Außenstehende 
diese nicht nur sehen, sondern selbst erleben?}

Im nächsten Schritt wurden konkrete Alltagssituationen analysiert. Durch Recherche und Gespräche 
wurde deutlich, dass insbesondere folgende Hindernisse immer wieder problematisch sind:

\begin{itemize}
    \item zu hoch angebrachte Bankautomaten
    \item nicht abgesenkte oder zu hohe Bordsteinkanten
    \item unebene Pflastersteine
\end{itemize}

Diese Beobachtungen führten zur Entscheidung, eine virtuelle Stadtumgebung zu entwickeln, in der genau 
diese Herausforderungen gezielt eingebaut werden.

Im Verlauf der Projektentwicklung wurde deutlich, dass eine rein visuelle VR-Simulation zwar die 
Perspektive verändern kann, jedoch das körperliche Empfinden nur eingeschränkt vermittelt. Aus diesem Grund 
entstand die weiterführende Idee, zusätzlich einen Hardware-Simulator zu entwickeln.

Ziel dieses Simulators war es, die Bewegungen aus der realen Welt präzise in die virtuelle Umgebung zu 
übertragen und die Bewegungen der VR-Welt physisch spürbar zu machen. Dadurch sollte nicht nur es nicht nur
visuell dargestellt werden, sondern auch ein mechanisches Feedback erzeugt werden.

Der entwickelte Aufbau ermöglicht es, reale Bewegungen, beispielsweise das Überfahren einer Kante oder einer 
unebenen Oberfläche, direkt an die Hardware-Simulation weiterzugeben. Gleichzeitig werden Bewegungsimpulse aus der 
Simulation mechanisch umgesetzt, sodass Nutzer*innen Höhenunterschiede oder Neigungen tatsächlich wahrnehmen 
können.